% GENERATED FILE. Edit canonical lessons, then run npm run generate:books.
% canonical-source-hash: fnv1a64:0be3ac02d4ad7827
% canonical-lessons: ES-C52-se-venden

\chapter{Se Compran Libros --- The Verb Tells On Itself}
\label{ch:se-agreement}
\hlchaptermodality{159}

\begin{chapteropening}
\textbf{By the end of this chapter:} \emph{I can make the verb agree with the thing, and say why it does.}

Complete ES-C52-se-venden: produce se compran libros and read it as if se still meant themselves.
\end{chapteropening}

\section[se compran libros]{se compran libros — the verb tells on itself}

\label{lesson:ES-C52-se-venden}



\begin{quote}
\emph{Se habla español} — one thing, singular verb. Now guess what happens with more than one thing.
\end{quote}

\begin{grammarlens}[title={Grammar Lens: it agrees, and that is strange}]
\begin{quote}
\emph{Se \textbf{compra} el libro.} — The book is bought. \emph{Se \textbf{compran} libros.} — Books are bought.
\end{quote}

The verb changed. \emph{Compra} for one book, \emph{compran} for several — the verb is agreeing with \textbf{the thing being bought}.

Stop and notice how odd that is. In \emph{se habla español} there is nobody doing the speaking, so what is the verb agreeing \textbf{with}? If this were simply \textquotedblleft{}one buys books,\textquotedblright{} the verb would stay singular, because \emph{one} is one person. It does not stay singular.

\textbf{The thing is behaving like the subject.} And that is the clue to where this construction came from.
\end{grammarlens}

\begin{cousinweb}
Read \emph{se compran libros} with the \emph{se} meaning what it meant in \emph{se llama}:

\begin{quote}
\textbf{Books buy themselves.}
\end{quote}

That is nonsense as a description of the world, and it is exactly the sentence Spanish is using. The reflexive was stretched: first \emph{the door closes itself} (true enough of doors), then \emph{the books sell themselves}, and finally any verb at all with the doer quietly dropped.

Once you see that, the agreement stops being strange. \emph{Libros} is the subject of \emph{se compran} in the same way \emph{él} is the subject of \emph{se llama}. The verb has been agreeing with its subject all along; it is just that the subject is a thing, and the thing is not really doing anything.

So this \emph{se} \textbf{is} the reflexive \emph{se} — grown, not borrowed. Unlike the \emph{se} of \emph{se lo digo}, which only looks like it.
\end{cousinweb}

\subsection*{Your turn}
\begin{itemize}
\raggedright
  \item \emph{Say it:} \textquotedblleft{}Se compra el libro\textquotedblright{} then \textquotedblleft{}Se compran libros\textquotedblright{} — hear the verb follow the thing
\end{itemize}

\begin{tcolorbox}[breakable,colback=teal!4,colframe=teal!35!black,title={Before you move on}]
Why does the verb go plural? (Because it agrees with \textbf{the thing}.) And why does it agree with the thing? (Because underneath, the thing is \textbf{doing it to itself}.) Next: the one you will use every day.
\end{tcolorbox}
